\documentclass[10pt,a4paper,oneside]{scrartcl}
\usepackage[english]{babel}
\usepackage{csquotes}
% Explicit light/dark entry points override this default.
\providecommand{\handouttheme}{light}
\expandafter\PassOptionsToPackage\expandafter{\handouttheme}{everett}
\usepackage[handout]{everett}
\usepackage{macro}
% An original Tufte-inspired layout on KOMA-Script and everett.sty.
% Reference: https://github.com/Tufte-LaTeX/tufte-latex/blob/master/sample-handout.tex
% Uses the reference's A4 column proportions; no Tufte class is required.
\usepackage{mathpazo}
\usepackage{tgpagella}
\usepackage[scale=0.92]{tgheros}
\usepackage[scaled=0.95]{inconsolata}
\usepackage{ragged2e}
\usepackage[oneside]{sidenotes}
\usepackage{marginfix}
% Footnotes are redirected below, so microtype's footnote patch is unnecessary.
\microtypesetup{nopatch=footnote}

\geometry{a4paper,left=24.8mm,textwidth=107mm,top=27.4mm,bottom=25mm,
  marginparsep=8.2mm,marginparwidth=49.4mm,
  headheight=16pt,headsep=8mm,footskip=12mm}
\normalmarginpar
\setlength{\marginparpush}{10pt}
\setstretch{1.08}
\setlength{\RaggedRightParindent}{1em}
\setlength{\emergencystretch}{2em}
\AtBeginDocument{\RaggedRight}
\setlist{topsep=5pt,itemsep=2pt,parsep=0pt,leftmargin=*}
\mdfsetup{startinnercode={\RaggedRight\setlength{\parindent}{0pt}},
  needspace=4\baselineskip}
\widowpenalty=10000
\clubpenalty=10000

% Quiet, unnumbered italic headings, as in the sample handout.
% Set secnumdepth to 2 in preamble.tex if numbered sections are preferred.
\setcounter{secnumdepth}{-1}
\setkomafont{disposition}{\normalfont\rmfamily\itshape\color{ctpText}}
\setkomafont{section}{\normalfont\rmfamily\itshape\Large\color{ctpText}}
\setkomafont{subsection}{\normalfont\rmfamily\itshape\large\color{ctpText}}
\RedeclareSectionCommand[beforeskip=16pt,afterskip=6pt]{section}
\RedeclareSectionCommand[beforeskip=12pt,afterskip=4pt]{subsection}
\renewcommand{\thesubsection}{\thesection.\arabic{subsection}}
\AddToHook{cmd/section/before}{\Needspace{7\baselineskip}}
\renewcommand{\vocab}[1]{\textbf{\color{ctpMauve}#1}}

\newenvironment{fullwidth}
  {\begin{adjustwidth}{0pt}{-\dimexpr\marginparwidth+\marginparsep\relax}}
  {\end{adjustwidth}}
\newcommand{\newthought}[1]{\par\addvspace{\baselineskip}\noindent\textsc{#1}}

% Unnumbered margin notes: \marginnote{...}; numbered notes: \sidenote{...}.
% Floating notes avoid one another; explicit offsets use fixed marginnote
% placement instead. Inspect dense margins after editing.
\newcommand{\handoutmarginfont}{\normalfont\footnotesize\color{ctpSubtext1}
  \setstretch{1.0}\RaggedRight\setlength{\parindent}{0pt}}
\makeatletter
\appto\@marginparreset{\handoutmarginfont}
\makeatother
\renewcommand*{\marginfont}{\handoutmarginfont}
\captionsetup{font={footnotesize,color=ctpSubtext1},labelfont=normalfont,
  justification=RaggedRight,singlelinecheck=false}
% sidenotes applies named styles locally; keep their captions ragged too.
\DeclareCaptionStyle{marginfigure}{font={footnotesize,color=ctpSubtext1},
  justification=RaggedRight,singlelinecheck=false}
\DeclareCaptionStyle{margintable}{font={footnotesize,color=ctpSubtext1},
  justification=RaggedRight,singlelinecheck=false}
\DeclareCaptionStyle{widefigure}{font={footnotesize,color=ctpSubtext1},
  justification=RaggedRight,singlelinecheck=false}
\DeclareCaptionStyle{widetable}{font={footnotesize,color=ctpSubtext1},
  justification=RaggedRight,singlelinecheck=false}
\DeclareCaptionStyle{sidecaption}{font={footnotesize,color=ctpSubtext1},
  justification=RaggedRight,singlelinecheck=false}
% Match the sample's footnotes-as-sidenotes behavior in the document body.
\AtBeginDocument{\let\footnote\sidenote}
% Keep everett's BibLaTeX/Biber setup. Ordinary \cite stays inline;
% \sidecite prints a full reference in a numbered margin note.
\newcommand{\sidecite}[1]{\sidenote{\fullcite{#1}.}}

\newcommand{\handoutcourse}[1]{\gdef\handoutcoursename{#1}}
\newcommand{\handoutrunningtitle}[1]{\gdef\handoutshorttitle{#1}}
\handoutcourse{}
\handoutrunningtitle{}
\KOMAoptions{headsepline=false,headwidth=textwithmarginpar,
  footwidth=textwithmarginpar}
\clearpairofpagestyles
\ihead{\handoutshorttitle}
\ohead{\pagemark}
\ofoot[\pagemark]{}
\setkomafont{pageheadfoot}{\normalfont\footnotesize\color{ctpSubtext1}}
\setkomafont{pagenumber}{\normalfont\footnotesize\color{ctpSubtext1}}

\makeatletter
\renewcommand{\maketitle}{%
  \par\begingroup
  \begin{fullwidth}
    \setlength{\parindent}{0pt}
    {\small\scshape\color{ctpSubtext1}\handoutcoursename\par}
    \vspace{5mm}
    {\rmfamily\itshape\fontsize{28}{32}\selectfont\@title\par}
    \ifx\@subtitle\@empty\else
      \vspace{2mm}{\large\itshape\color{ctpSubtext1}\@subtitle\par}
    \fi
    \vspace{5mm}
    {\small\@author\par}
    {\small\color{ctpSubtext1}\@date\par}
    \@thanks
  \end{fullwidth}
  \endgroup\thispagestyle{plain}\vspace{3mm}}
\makeatother
\renewenvironment{abstract}
  {\par\small\noindent\ignorespaces}
  {\par\vspace{4pt}}

% Theorem counters run through the handout, independently of sections.
\declaretheoremstyle[headfont=\normalfont\itshape\color{ctpMauve},
  bodyfont=\normalfont,spaceabove=8pt,spacebelow=4pt]{handoutexercise}
\declaretheorem[style=handoutexercise,name=Exercise]{exercise}
\crefname{exercise}{exercise}{exercises}
\Crefname{exercise}{Exercise}{Exercises}

% Theme-aware defaults for figures supplied by the author.
\tikzset{every picture/.append style={color=ctpText}}
\tikzcdset{every diagram/.append style={color=ctpText},
  arrows={color=ctpText},labels={text=ctpText},background color=ctpBase}
\pgfplotsset{every axis/.append style={
  axis line style={draw=ctpSubtext1},tick style={draw=ctpSubtext1},
  tick label style={color=ctpText},label style={color=ctpText},
  legend style={fill=ctpMantle,draw=ctpSurface1,text=ctpText},
  grid style={draw=ctpSurface0}}}


% Put project-specific settings and notation here.
\crefname{definition}{definition}{definitions}
\Crefname{definition}{Definition}{Definitions}
\crefname{prop}{proposition}{propositions}
\Crefname{prop}{Proposition}{Propositions}


% Edit these fields and replace the sample material below.
\title{Groups and Symmetry}
\subtitle{A mathematical handout}
\author{Everett}
\date{\today}
\handoutcourse{Abstract Algebra / Handout 01}
\handoutrunningtitle{Groups and symmetry}

\begin{document}
\maketitle

\begin{abstract}
A short introduction to groups, illustrated by the symmetries of a square.
We connect the axioms to a concrete picture, prove cancellation, and end
with exercises. This source also demonstrates the handout's margin notes,
figures, citations, and full-width material.
\end{abstract}

\section{From a picture to an operation}

\newthought{Start with a square.} Rotate it through a quarter turn, then
rotate it again. The result is another symmetry: a half turn.
The operation is composition, and its inputs and outputs all belong to
the same collection.\sidenote{Read a composition from right to left:
$gh$ means apply $h$ first, then $g$. Fix this convention before computing.}

\begin{marginfigure}
\centering
\begin{tikzpicture}[scale=0.65, color=ctpText]
  \draw[thick,ctpBlue,fill=ctpBlue!6!ctpBase] (0,0) rectangle (2,2);
  \foreach \x/\y/\name in {0/0/D,2/0/C,2/2/B,0/2/A}
    \fill (\x,\y) circle (2pt) node[shift={({(\x-1)*0.25},{(\y-1)*0.25})}] {$\name$};
  \draw[-{Latex},ctpTeal,thick] (2.7,0.5) arc[start angle=-35,end angle=70,radius=1];
  \node at (3.0,1.4) {$r$};
\end{tikzpicture}
\caption{A quarter turn $r$ cycles the four vertices. Four quarter turns
return every vertex to its starting position.}
\label{fig:square}
\end{marginfigure}

\begin{definition}[Group]\label{def:group}
A \vocab{group} is a set $G$ with an associative operation
$G\times G\to G$, an identity $e$, and an inverse $g\inv$ for each $g\in G$.
Thus $eg=ge=g$ and $gg\inv=g\inv g=e$.
\end{definition}

\begin{eg}
The integers $\ZZ$ form a group under addition. The identity is $0$,
and the inverse of $n$ is $-n$. The rotations of the square form a group
$\{e,r,r^2,r^3\}$, with $r^4=e$.
\end{eg}

\subsection{The order of an operation matters}
Rotating and reflecting a square need not commute. If $s$ is a reflection,
then $sr=r\inv s$. A group is \vocab{abelian} when $gh=hg$ for all
$g,h\in G$. The rotation group is abelian; the full symmetry group is not.
A standard reference develops these examples further
\sidecite{judson2020}.

\section{Cancellation comes from inverses}

\begin{theorem}[Left cancellation]\label{thm:cancellation}
If $a,b,c\in G$ and $ab=ac$, then $b=c$.
\end{theorem}
\begin{proof}
Multiplying on the left by $a\inv$ gives
\[
  b=(a\inv a)b=a\inv(ab)=a\inv(ac)=(a\inv a)c=c.
\]
Associativity allows us to regroup each product.
\end{proof}

The argument uses the inverse axiom in \cref{def:group}.
It never changes the order of the factors, so
\cref{thm:cancellation} also applies to nonabelian groups.

% A deliberate page break keeps this two-page teaching example easy to edit.
% Remove it to let a longer handout flow naturally.
\clearpage
\section{Read a finite group from its table}

\newthought{A table records every product.} For the rotation group, the
entry in row $r^i$ and column $r^j$ is $r^{i+j}$, reducing the exponent
modulo $4$. Each row and each column contains every element exactly once.

\begin{margintable}
\centering
\begin{tabular}{c|cccc}
  \toprule
  $\cdot$ & $e$ & $r$ & $r^2$ & $r^3$ \\
  \midrule
  $e$   & $e$   & $r$   & $r^2$ & $r^3$ \\
  $r$   & $r$   & $r^2$ & $r^3$ & $e$   \\
  $r^2$ & $r^2$ & $r^3$ & $e$   & $r$   \\
  $r^3$ & $r^3$ & $e$   & $r$   & $r^2$ \\
  \bottomrule
\end{tabular}
\caption{The multiplication table of the four rotations.}
\label{tab:rotations}
\end{margintable}

\begin{remark}
A repeated entry in a row would contradict left cancellation.
Right cancellation gives the corresponding statement for columns.
\end{remark}

\subsection{A visual summary}
\Cref{fig:cycle} uses the main text column and the note column together.
Use \texttt{figure*} or \texttt{table*} for wide floats; use
\texttt{fullwidth} for a wide passage of text. Keep margin notes away
from the vertical space occupied by wide content.

\begin{figure*}[htbp]
\centering
\begin{tikzpicture}[
  every node/.style={text=ctpText},
  state/.style={draw=ctpBlue,fill=ctpBlue!6!ctpBase,
    rounded corners=2pt,minimum width=23mm,minimum height=12mm},
  every path/.style={draw=ctpTeal,thick,-{Latex}}
]
  \node[state] (e) {$e$};
  \node[state,right=17mm of e] (r) {$r$};
  \node[state,right=17mm of r] (rr) {$r^2$};
  \node[state,right=17mm of rr] (rrr) {$r^3$};
  \draw (e) -- node[above] {$r$} (r);
  \draw (r) -- node[above] {$r$} (rr);
  \draw (rr) -- node[above] {$r$} (rrr);
  \draw (rrr.south) -- ++(0,-0.55) -| node[pos=0.25,below] {$r$} (e.south);
\end{tikzpicture}
\caption{Successive quarter turns visit each rotation and return to the identity.}
\label{fig:cycle}
\end{figure*}

\section{Try it yourself}
\begin{exercise}\label{ex:right}
Prove right cancellation: if $ba=ca$, then $b=c$.
Identify exactly where you use associativity.
\end{exercise}

\begin{exercise}
Show that every element has a unique inverse. Then find the inverse
of each rotation in \cref{tab:rotations}.
\end{exercise}

\begin{exercise}
Let $s$ reflect the square across its vertical axis. Track a vertex to
show that $rs\ne sr$.
\end{exercise}

\begin{fullwidth}
\small\itshape
Takeaway: a symmetry is an operation that preserves a structure.
The group axioms capture composition, doing nothing, and undoing an
operation. Cancellation follows from these axioms rather than from
commutativity.
\end{fullwidth}

\printbibliography[heading=bibintoc,title={Further reading}]
\end{document}
